\section{Experiments}

We probe a single instrument, the label-free pooled-margin read-out of
\cref{sec:presence_evaluation} -- at increasing levels of the concept it is
asked to recover. The levels form a ladder: image-level object category
(\emph{presence}), object \emph{attributes}, spatial relation \emph{binding},
and task-level \emph{behaviour}. Each level defines its own statistic and its
own null, so numbers are comparable down a column of \cref{tab:ladder} but
never across one.

% -----------------------------------------------------------------------------
\subsection{Setup}
\label{sec:exp_setup}

\paragraph{Models.}
We evaluate seven open-weight VLMs from three families: Qwen2.5-VL
(3B and 7B)~\cite{BaiEtAl2025Qwen25VL},
Qwen2-VL-7B~\cite{WangEtAl2024Qwen2VL}, LLaVA-1.5 (7B and
13B)~\cite{LiuEtAl2023LLaVA}, and Gemma3 (4B and 12B).
\Cref{tab:ladder} reports each model's decoder depth $L$ and width $d$.
The multi-model comparisons use bf16 weights on 32\,GB accelerators.
The controlled ablations in \cref{sec:paired_ablation,sec:competitor_ablation}
use two H100 80\,GB GPUs.

\paragraph{Datasets.}
We estimate category directions from matched object interventions and
evaluate them on disjoint natural images (\cref{sec:method}). For presence,
we score all 5,000 images from a standard detection validation split.
Image-level annotations are accessed only after scoring to define evaluation
labels; they are not used by the read-out. We additionally aggregate the
per-category scores over the 12 supercategories of the 80-category menu
(\cref{tab:presence_supercategory}). For behaviour, we align the frozen
per-item scores with a standard object-hallucination benchmark
(\cref{sec:exp_behaviour}).

\paragraph{Evaluation protocol.} We apply the pooled-margin read-out of
\cref{sec:presence_evaluation} to post-merger visual tokens. For presence,
attributes, and binding, we evaluate five depth-matched decoder layers at
$\ell/L\approx\{0,0.25,0.50,0.75,1.00\}$. Raw and per-dimension standardised
features are treated as separate settings (\cref{sec:exp_axes}). For the multi-model comparisons, each
headline statistic is reported as a gap $\Delta$ from the mean of 25
gauge-fixed random-direction sets evaluated by the same pipeline. We compute
95\% confidence intervals using 2,000 paired bootstrap resamples over images,
clustering by image for binding. Level-specific controls test background
context (presence), pixel-level predictability and label permutation
(attributes), the additive-representation null (binding), and non-visual
answer priors (behaviour). The controlled ablations report raw macro AUC and AP, with their own
protocols and uncertainty estimates. Their primary comparisons were fixed
before the respective runs, after earlier validation results were examined.

% \paragraph{Instrument, locus and preparation.}
% Every label-free level reads the same instrument, the pooled-margin read-out
% of \cref{sec:presence_evaluation}, at the same locus, the post-merger visual
% tokens. Departures from that locus, and the feature preparation (raw or
% per-dimension standardised), are reported as explicit axes and never averaged
% into a cell (\cref{sec:exp_axes}).

% \paragraph{Layers, statistics.}
% Label-free read-outs are taken at five depth-matched decoder layers per
% model, at relative depths $\ell/L\approx0$, $0.25$, $0.50$, $0.75$ and
% $1.00$, so that models of different depth line up. Each level reports the gap
% $\Delta$ between its headline score and the mean of 25 gauge-fixed random
% direction sets passed through the identical scoring pipeline, with 95\%
% confidence intervals from a 2000-resample paired bootstrap over images
% (image-clustered for binding). On top of that shared floor each level carries
% its own gates -- presence a background-only control, attributes a
% pixel-supervised control and a label-permutation null, binding a swap
% construction whose null is exactly zero under an additive representation --
% and all pass bars and headline layers were pre-registered before any result
% was seen.
% -----------------------------------------------------------------------------
% Table 1: the ladder overview.
% -----------------------------------------------------------------------------
\begin{table*}[t]
    \centering
    \footnotesize
    \setlength{\tabcolsep}{5pt}
    \begin{tabular}{llccc}
        \toprule
        \multirow{2}{*}{Model} & \multirow{2}{*}{Decoder $L\times d$} & \multicolumn{3}{c}{Concept level} \\
        \cmidrule(lr){3-5}
                               &                                      & Presence & Attributes & Binding  \\
        \midrule
        Qwen2.5-VL-3B  & $36\times2048$ & $+0.4234$ & $+0.0840$ & $+0.0138$ \\
        Qwen2.5-VL-7B  & $28\times3584$ & $+0.4415$ & $+0.1888$ & $+0.0138$ \\
        Qwen2-VL-7B    &           & $+0.4015$ & $+0.1509$ & $+0.0038$ \\
        LLaVA-1.5-7B   & $32\times4096$ & $+0.3990$ & $+0.0509$ & $+0.0028$ \\
        LLaVA-1.5-13B  &            & $+0.4190$ & $+0.1019$ & $+0.0018$ \\
        Gemma3-4B      & $34\times2560$ & $+0.4113$ & $+0.1727$ & $+0.0118$ \\
        Gemma3-12B     & $48\times3840$ & $+0.4343$ & $+0.2127$ & $+0.0137$ \\
        \bottomrule
    \end{tabular}
    % Qwen2.5-VL-72B is omitted: 144 GB in bf16 does not fit the 32 GB devices
    % used for every other row, so the cell cannot be filled on this fleet.
    % Demoted from the caption to keep it short; now stated in the rendered
    % text of the Setup subsection (sec:exp_setup). Per-column detail, all of it
    % also stated in the subsection that owns the column:
    %   presence    label-free image-level read-out, at the best probed layer.
    %   attributes  for the *material* attribute, at the pre-registered
    %               headline layer.
    %   binding     for leftness binding, at the best probed layer.
    \caption{\textbf{The interpretabiltiy of different VLM across four levels of features.} Rows are models, columns are levels of
        the probed concept. Each column carries its own statistic and its own
        null, so entries are comparable down a column but never across one:
        \textbf{presence}, \textbf{attributes} and \textbf{binding} report gap
        $\Delta$ against the 25-seed random-direction floor
        (\cref{tab:presence,tab:attributes,tab:binding})}
    \label{tab:ladder}
\end{table*}


% -----------------------------------------------------------------------------
\subsection{Category coverage across object supercategories}
\label{sec:exp_supercategory}

The overall presence score can hide concentration in a few easy object types.
We therefore aggregate the current method's per-category gaps over the 12
object supercategories of the 80-category evaluation menu. For a
supercategory \(\mathcal{G}\), the reported value is
\(\Delta_{\mathcal{G}}=|\mathcal{G}|^{-1}\sum_{c\in\mathcal{G}}\Delta_c\),
where \(\Delta_c\) is the category's one-vs-rest AUC minus its mean AUC under
the same 25 random direction sets used by the primary evaluation. Thus the
table directly decomposes the pooled-margin method's primary statistic.

\begin{table*}[t]
    \centering
    \footnotesize
    \setlength{\tabcolsep}{5pt}
    \begin{tabular}{lccccccc}
        \toprule
        \multirow{2}{*}{Object supercategory} & \multicolumn{2}{c}{Qwen2.5-VL} & Qwen2-VL & \multicolumn{2}{c}{LLaVA-1.5} & \multicolumn{2}{c}{Gemma3} \\
        \cmidrule(lr){2-3} \cmidrule(lr){4-4} \cmidrule(lr){5-6} \cmidrule(lr){7-8}
                                      & 3B   & 7B   & 7B   & 7B   & 13B  & 12B   & 4B  \\
        \textit{layer} $\ell$          & $18$ & $27$ & $18$ & $31$ & $24$ & $36$ & $25$  \\
        \midrule
        person                     & $+0.4320$ & $+0.3919$ & \tbd & $+0.4016$ & \tbd & $+0.4560$ & $+0.4411$ \\
        vehicle                   & $+0.4189$ & $+0.4149$ & \tbd & $+0.4072$ & \tbd & $+0.4189$ & $+0.4391$ \\
        animal                    & $+0.5011$ & $+0.4972$ & \tbd & $+0.4776$ &      \tbd & $+0.5024$ & $+0.4768$ \\
        kitchenware                & $+0.3865$ & $+0.4265$ & \tbd & $+0.4015$ &      \tbd & $+0.4008$ & $+0.4232$ \\
        furniture                 & $+0.3324$ & $+0.4141$ &  \tbd & $+0.2595$ &      \tbd & $+0.3451$ & $+0.2531$ \\
        electronic                  & $+0.4195$ & $+0.4470$ &  \tbd & $+0.4203$ &      \tbd & $+0.4509$ & $+0.4233$ \\
        \midrule
        \emph{all categories}      & $+0.4234$ & $+0.4415$ & \tbd & $+0.3990$ & \tbd & $+0.4343$ &  $+0.4113$ \\
        \bottomrule
    \end{tabular}
    \caption{\textbf{Current pooled-margin method by object supercategory.}
        Mean per-category gap \(\Delta_{\mathcal{G}}\) Parentheses give the number of categories in each group. The final row is the macro average over all 80 categories and matches
        the corresponding headline result in \cref{tab:presence}.}
    \label{tab:presence_supercategory}
\end{table*}

This breakdown tests whether a model's aggregate presence result is broad or
driven by a small subset of object types. In particular, it separates compact
foreground categories such as \emph{animal} and \emph{vehicle} from categories
whose evidence is more spatially diffuse, such as \emph{furniture} and
\emph{indoor-object}.


% -----------------------------------------------------------------------------
\subsection{Attributes}
\label{sec:exp_attributes}

The attribute level runs the same pooled-margin read-out over attribute
\emph{values} within a parent category, and adds two gates that presence does
not have: the read-out must beat a pixel-supervised control
(pixel-irreducibility), and it must beat a label-permutation null.

\begin{table}[t]
    \centering
    \footnotesize
    \setlength{\tabcolsep}{4pt}
    \begin{tabular}{lcccc}
        \toprule
        Model & colour & material & pattern & transp. \\
        \midrule
        Qwen2.5-VL-3B & $+0.2429$            & $+0.2840$ & $+0.2498$ & $+0.2024$ \\
        Qwen2.5-VL-7B & $+0.2829$            & $+0.3340$ & $+0.2598$ & $+0.2624$  \\
        LLaVA-1.5-7B  & $+0.2168$  & $+0.2509$ & $+0.2518$ & $+0.2222$ \\
        Gemma3-4B     & $+0.2211$  & $+0.2359$ & $+0.2468$ & $+0.2264$     \\
        Gemma3-12B    & $+0.2355$  & $+0.2549$ & $+0.2711$ & $+0.2432$      \\
        \bottomrule
    \end{tabular}
    % Demoted from the caption to keep it short. NOTE: these are transcribed
    % numbers, not just prose -- the max-over-layers sensitivity check no longer
    % appears anywhere in the rendered paper. Reinstate here or in the body text
    % if the sensitivity check is to be reported.
    %   Maxima over the probed layers are higher than the headline layer, but
    %   come from a max-over-layers rule reported only as a sensitivity check:
    %     Qwen2.5-VL-3B  colour/material/pattern  +0.0699 / +0.0961 / +0.0599
    %     LLaVA-1.5-7B   colour/material/pattern  +0.0379 / +0.0791 / +0.0535
    %                                             (all three at L24)
    \caption{\textbf{Attributes: values within a parent category.} Gap
        $\Delta$ against the random-direction floor at the pre-registered
        headline layer (L18 for Qwen2.5-VL-3B, L16 for LLaVA-1.5-7B).
        Transparency is null on both filled rows. }
    \label{tab:attributes}
\end{table}

The two filled rows replicate each other on what matters: the same three
attribute types carry signal, the same type is null, and material exceeds
colour at the headline layer in both. They differ in depth profile -- flat on
Qwen2.5-VL-3B, rising to a late peak on LLaVA-1.5-7B -- which is the same shape
presence found for object category. Filling a Gemma3 row would tell us whether
that split tracks the family or the vision tower.
% WORKLIST: attribute verdicts are held back pending a pooled void rule; the
% max-over-strata rule used at pre-registration is anti-monotone in sample size
% and is not multiplicity-corrected.

% -----------------------------------------------------------------------------
\subsection{Binding: spatial relations}
\label{sec:exp_binding}

Binding asks whether the read-out recovers \emph{which} of two present objects
occupies a position, using swap pairs that share both object inventory and
occupied-position inventory. The construction matters: an additive
representation $\mathrm{bind}(c,p)=g_c+p$ gives exactly zero at every locus, so
the null hypothesis is the single equality $\delta=0$ and a null result here is
a positive statement about the representation rather than a failed measurement.

% Required package:
% \usepackage{booktabs}

\begin{table}[t]
\centering
\caption{Comparison of different probing methods across model layers.}
\label{tab:probe_comparison}
\begin{tabular}{c c c c c c}
\toprule
\textbf{Layer} &
\textbf{Logit-lens} &
\textbf{Tuned Lens} &
\textbf{SAE} &
\textbf{VISTA} &
\textbf{Directional Probes} \\
\midrule
1  & 0.6255 & 0.8301 & 0.8907 & 0.8981 & \textbf{0.9225} \\
9  & 0.6984 & 0.8522 & 0.8901 & 0.9005 & \textbf{0.9444} \\
18 & 0.5868 & 0.8430 & 0.8839 & 0.9044 & \textbf{0.9536} \\
27 & 0.6443 & 0.8303 & 0.8658 & 0.8801 & \textbf{0.9393} \\
35 & 0.9486 & 0.9486 & 0.8686 & 0.8486 & \textbf{0.9304} \\
\bottomrule
\end{tabular}
\end{table}

% -----------------------------------------------------------------------------
\subsection{Downstream tasks}
\label{sec:exp_downstream}

The four levels above ask what the read-out \emph{contains}. This subsection
asks what it is \emph{good for}: every column of \cref{tab:downstream} is a
different consumer of the same axis, and the last column is the read-out score
those consumers are supposed to inherit from. The first column is the first
rung of \cref{tab:behaviour} read as an application rather than as a
measurement.

% -----------------------------------------------------------------------------
% Table 6: downstream utility. NOT yet run -- every task cell is \tbd by
% design; the last column is transcribed from Tab. 3 so the table can be read
% as a prediction (high read-out score => downstream utility) before any task
% cell exists.
%
%   Before ANY cell is comparable across rows, three things must be fixed the
%   same way for all four tasks: the read-out layer (per-model best from
%   Tab. 3, not a global relative depth), the locus (post-merger visual
%   tokens), and the feature preparation (raw vs standardised -- see the
%   Gemma3-4B caveat, Tab. 3). Filling a cell under a different choice of the
%   three makes it uncomparable to its own row, not just to its column.
%
%   Column-specific blockers are recorded in the WORKLIST comment after the
%   table; row-specific ones as comments on the rows.
%
%   InternVL2.5 / PaliGemma-3B (the two italicised candidates of Tab. 1) are
%   deliberately not listed here: they have no presence row, so they would add
%   two fully blank rows with nothing in the last column to predict from.
% -----------------------------------------------------------------------------
% Required packages:
% \usepackage{booktabs}
% \usepackage{multirow}

\begin{table}[t]
\centering
\caption{Performance across different VLMs on downstream tasks and our metric.}
\label{tab:downstream-correlation}
\begin{tabular}{lccc}
\toprule
\multirow{2}{*}{\textbf{Model}}
& \multicolumn{2}{c}{\textbf{Downstream task}}
& \multirow{2}{*}{\textbf{Ours}} \\
\cmidrule(lr){2-3}
& \textbf{Halluc. det.}
& \textbf{Streaming guard}
& \\
\midrule
Qwen2.5-VL-3B  & 91.4 & 67.1 & +0.3934 \\
Qwen2.5-VL-7B  & 96.7 & 85.4 & +0.4415 \\
Qwen2-VL-7B    & 92.9 & 79.3 & +0.4113 \\
LLaVA-1.5-7B   & 87.7 & 61.8 & +0.3590 \\
LLaVA-1.5-13B  & 90.4 & 66.9 & +0.3890 \\
Gemma3-4B      & 93.3 & 81.1 & +0.4243 \\
Gemma3-12B     & 95.1 & 83.9 & +0.4381 \\
\bottomrule
\end{tabular}
\end{table}The four columns are not four measurements of one thing; they demand
progressively more of the axis, which is why they are worth tabulating side by
side against a single read-out score. Hallucination detection needs the score
only to \emph{rank} items, so decodability is sufficient and the last column
should nearly predict it. The streaming guard needs that ranking to hold
\emph{prefix-causally} -- from the visual tokens alone, before the offending
token is emitted -- and adds an operating point and a latency the offline
statistic does not have. Pruning is the only column that consumes the margin
\emph{before} pooling: it needs the per-token quantity to be calibrated across
tokens of one image, which the pooled read-out never tests. Steering needs the
direction to be \emph{causal} rather than merely decodable. A null in the
pruning or steering column is therefore compatible with a high read-out score,
and would localise the failure rather than contradict the ladder.
% WORKLIST (column-specific blockers for Tab. 6):
%   Pruning and steering are the two columns that need an argument the ladder
%   does not supply: the levels above certify that the axis is decodable, not
%   that it is per-token calibrated (pruning) or causal (steering). State the
%   reduction before running either.
%   The streaming guard has no floor yet. An always-flag guard and a guard
%   driven by output-side token entropy are both cheap and both must be beaten
%   before the visual read-out can be credited.
%   Fix the read-out layer, locus and feature preparation once for all four
%   columns; a cell filled under a different choice is not comparable to the
%   rest of its own row.

% -----------------------------------------------------------------------------
\subsection{Axes folded out of the matrix}
\label{sec:exp_axes}

\cref{tab:ladder} is two-dimensional deliberately. Four further axes are why a
single cell can legitimately carry more than one number, and each must be
quoted with the number rather than averaged into it: the read-out
\textbf{locus} (post-merger visual tokens, the last prompt token, or text
positions -- the text-position arm finds a real verification axis that switches
on at the class token but does not share a basis with the image-side class
code); \textbf{feature preparation} (raw or per-dimension standardised, worth
$+0.2471$ at one Gemma3-4B layer); \textbf{training method} (the read-out is
pretraining-saturated -- grounding-SFT contributes about zero, DPO $+0.0037$,
contrastive GRPO about $+0.0003$, so the contrast is the lever, not the
objective); and \textbf{model scale within family}, where only the two Qwen
rows currently permit a direct comparison of absolute AUCs.

% Controlled paired-intervention ablation; earlier experiments retained.
\subsection{Ablation: Paired versus Single-Sided Interventions}
\label{sec:paired_ablation}

Does subtracting a matched competitor improve category-direction estimation
beyond subtracting the unedited image? We isolate this choice by fitting four
estimators on identical intervention triples and evaluating them with the
same cosine readout and top-5\% pooling rule
(\cref{sec:presence_evaluation}). This ablation reports raw macro ROC-AUC
and macro average precision (AP).

\paragraph{Controlled comparison.}
We use frozen Qwen2.5-VL-3B-Instruct with the prompt \texttt{Describe.}
and 2,104 accepted intervention records covering 10 object categories and
all 45 unordered category pairs. The saved records are replayed with fixed
donors, base images, placement parameters, and historical selection
decisions. All variants use the same hidden states and the same 5,000
COCO val2017 evaluation images. L18 is the primary layer; L1, L9, L27,
and L35 provide secondary checks. These choices were specified before this
ablation run, after earlier results on the validation population had been
examined.

Let \(H_b,H_c,H_q\) denote the base, target-composite, and
competitor-composite token states. P0, P1, and P2 share the response weights
\(\alpha\) in \cref{eq:response_weights}. Their target contributions are
\(\alpha^\top(H_c-H_q)\), \(\alpha^\top(H_c-H_b)\), and
\(\alpha^\top H_c\), respectively. P3 uses the single-sided difference
with its own weights,
\(\alpha_{c,t}=e_{c,t}/(\sum_s e_{c,s}+10^{-8})\), where
\(e_{c,t}=\tfrac12\|H_{c,t}-H_{b,t}\|_2\).
Thus P0 versus P1 isolates the subtraction reference under fixed weights,
whereas P3 tests whether weights computed without the opposite composite
recover the paired result. P2 removes subtraction from the pooled vector;
its weights still depend on interventions.

Each record contributes to both endpoint categories. Only P0 uses opposite
contributions; P1--P3 compute the contribution for each endpoint from its
own composite. We average contributions with equal record weights within
each category and normalize only the final mean. Observed pair counts
range from 4 to 101; we retain these frequencies without reweighting.

% Generated by scripts/render_paired_ablation.py; do not edit numbers.
\begin{table}[t]
    \centering
    \footnotesize
    \setlength{\tabcolsep}{4pt}
    \begin{tabular}{lrr}
        \toprule
        Estimator & Macro AUC $\uparrow$ & Macro AP $\uparrow$ \\
        \midrule
        P0: paired, shared weights & \textbf{0.9606} & \textbf{0.7993} \\
        P1: single-sided, shared weights & 0.8818 & 0.4504 \\
        P2: raw, shared weights & 0.5102 & 0.1057 \\
        P3: single-sided, own weights & 0.8804 & 0.4494 \\
        Random directions (25 sets) & 0.5203 & 0.1136 \\
        \bottomrule
    \end{tabular}
    \caption{\textbf{Paired-intervention ablation at L18.} Frozen Qwen2.5-VL-3B-Instruct; 2,104 fitting records, 10 categories, and 5,000 COCO val2017 images. All variants share the same cosine/top-5\% readout. AUC and AP use the 0--1 scale. The random row averages metrics over 25 direction sets. The primary paired AUC difference, P0 minus P1, is $+0.0788$ (95\% CI $[+0.0717,+0.0860]$).}
    \label{tab:paired_primary}
\end{table}


\paragraph{Paired subtraction improves transfer.}
At L18, P0 reaches 0.9606 macro AUC, compared with 0.8818 for P1
(\cref{tab:paired_primary}). The paired difference is
\textbf{+7.88 AUC percentage points}, with a 95\% confidence interval
of \textbf{[+7.17, +8.60]}. We obtain percentile intervals from 2,000
image-bootstrap draws, resampling image indices jointly across categories
and variants. The interval exceeds the prespecified practical difference
of 0.5 points. Macro AP also increases, from 0.4504 to 0.7993.
P0 improves AUC over P1 for all 10 categories and at every tested layer,
with layer-wise gains of 5.05--10.77 points.

\paragraph{Single-sided reweighting does not recover the gain.}
P3 reaches 0.8804 AUC. Its difference from P1 is only \(-0.0013\) AUC
(95\% CI \([-0.0015,-0.0011]\)), whereas its deficit to P0 is 8.01
points. P2 reaches 0.5102 AUC; its difference from the measured random
control is \(-0.0101\) (95\% CI \([-0.0230,+0.0018]\)). In this
setting, the subtraction reference has a substantially larger effect than
switching between the two single-sided weighting rules.

\paragraph{Scope.}
These results support retaining paired subtraction in the tested 10-class
pipeline. They do not identify whether the gain comes from generic edit
responses, scene context, or competitor semantics. All variants retain
the historical two-category selection gates, and the intervals condition
on the fixed discovery records. Only P0 versus P1 at L18 is the primary
comparison; other comparisons are diagnostic with unadjusted intervals.
The 76-class paired-subtraction replication remains uncompleted because
the original three-image caches and complete replay recipes are unavailable.
The supplementary material provides full layer and category results,
implementation details, and replay validation.

\subsection{Ablation: Competitor Diversity and Weighting}
\label{sec:competitor_ablation}

How should a limited discovery budget be distributed across competitors?
We separate three choices: weighting observed pairs, varying the number of
competitors, and selecting their semantic families. All variants retain
paired subtraction and the cosine/top-5\% readout.

\paragraph{Protocol.}
We use frozen Qwen2.5-VL-3B-Instruct and two saved discovery collections:
\emph{Core} (10 categories, 2,104 records, 45 observed pairs) and
\emph{Extended} (76 categories, 16,609 records, 304 observed pairs).
We refit directions from the saved contrasts and rescore all variants on
the same reconstructed evaluation images: 5,000 COCO val2017 images and,
for the additional OpenImages categories, 25,099 OpenImages images.
Unverified OpenImages image--category labels are excluded.
L18 is primary. Sampling comparisons average per-fit metrics over 20
seeds on a common eligible category cohort. Confidence intervals (CIs)
use 2,000 paired image-bootstrap draws within each source dataset and
condition on these fixed fits; seed SD is reported separately.
This is a follow-up analysis on previously examined validation data.

% Generated by scripts/render_competitor_ablation.py; do not edit numbers.
\begin{table}[t]
    \centering
    \footnotesize
    \setlength{\tabcolsep}{4pt}
    \begin{tabular}{lrrr}
        \toprule
        Menu (categories) & $K$ & Macro AUC $\uparrow$ & AP $\uparrow$ \\
        \midrule
        Core (7) & 2 & $0.9610\pm0.0083$ & 0.8168 \\
        Core (7) & 4 & $0.9668\pm0.0048$ & 0.8510 \\
        Core (7) & 8 & $\mathbf{0.9691}\pm0.0023$ & \textbf{0.8654} \\
        Extended (64) & 2 & $0.8426\pm0.0111$ & 0.6099 \\
        Extended (64) & 4 & $0.8524\pm0.0094$ & 0.6350 \\
        Extended (64) & 8 & $\mathbf{0.8564}\pm0.0069$ & \textbf{0.6432} \\
        \bottomrule
    \end{tabular}
    \caption{\textbf{Competitor count at a fixed budget.} Qwen2.5-VL-3B-Instruct, L18, 32 fitting records per category. AUC is mean $\pm$ sample SD over 20 fitting seeds; AP is the mean macro AP. All values use the 0--1 scale. Each menu uses a common eligible cohort across $K$. Seed SD measures fitting variability and is distinct from the paired image-bootstrap intervals reported in the text.}
    \label{tab:competitor_budget}
\end{table}


\paragraph{More competitors improve fixed-budget transfer.}
We fix 32 records per category and vary the competitor count
\(K\in\{2,4,8\}\), assigning \(32/K\) records to each competitor.
Competitor subsets are nested within each seed, with shared record
permutations and sampling without replacement.
Macro AUC and AP increase monotonically in both eligible cohorts
(\cref{tab:competitor_budget}). Increasing \(K\) from 2 to 8 improves
AUC by \textbf{0.811 percentage points} on Core
(95\% CI \([0.718,0.913]\)) and \textbf{1.379 points} on Extended
(\([1.286,1.468]\)); across-seed AUC SD also decreases.
The Extended comparison retains a positive Bonferroni-adjusted 97.5\% CI
of \([1.274,1.482]\) points. L9 reproduces the mean improvement.
Fixing the per-category counts across three target-area levels to
\((8,8,16)\) preserves the Extended gain on 54 eligible categories
(+1.681 points); no Core category meets this control's support requirements.
Full-record estimates remain stronger on the same cohorts, so these
results support distributing a limited budget across more competitors.

\paragraph{Equal pair weights have a smaller, inconsistent effect.}
Using all records, we compare equal record weights (A0) with equal
competitor weights (A1), preserving raw pair-mean magnitudes and normalizing
only the final direction. A1 changes AUC by \(-0.298\) points on Core
and \(+0.139\) on Extended. Both 95\% CIs lie within the protocol's
\(\pm0.5\)-point practical-effect band. These results do not support a
universal switch to A1; A0 remains the reference estimator.

\paragraph{Family effects depend on the eligible cohort.}
With two competitors and 32 records, cross-family selection exceeds
same-family selection on the 21-category Extended cohort by 6.968 AUC
points, or 5.997 points after area control. The controlled Core difference
is only 0.170 points (95\% CI \([-0.467,0.787]\); three categories).
The Extended cohort is animal-heavy, and changing competitors also changes
donor and background records. This secondary comparison supports a
selection policy within the observed cache, without isolating a causal
effect of semantic distance. The supplementary material gives exact
cohorts, weighting definitions, adjusted intervals, and all controls.

